%----------------------------------------------------------------------------------------
% Generic one-page Chinese CV template
%----------------------------------------------------------------------------------------
\documentclass[UTF8,a4paper,10pt]{ctexart}

\usepackage{parskip}
\usepackage[margin=0.55in]{geometry}
\usepackage{titlesec}
\usepackage[unicode,hidelinks]{hyperref}

\titleformat{\section}{\large\bfseries\raggedright}{}{0em}{}[\titlerule]
\titlespacing{\section}{0pt}{4pt}{2pt}
\setlength{\parindent}{0pt}
\hyphenpenalty=5000
\tolerance=1000
\emergencystretch=1em
\clubpenalty=10000
\widowpenalty=10000
\brokenpenalty=10000
\raggedbottom

\newcommand{\cvneedspace}[1]{%
  \par\ifdim\dimexpr\pagegoal-\pagetotal\relax<#1\relax\newpage\fi
}
\newcommand{\cvsection}[1]{\cvneedspace{4\baselineskip}\section{#1}}

\begin{document}
\pagestyle{empty}

\begin{center}
  {\LARGE\bfseries 候选人姓名} \\[4pt]
  {\small
    \href{https://github.com/your-handle}{github.com/your-handle} \ $|$ \
    \href{mailto:you@example.com}{you@example.com} \ $|$ \
    \href{https://www.linkedin.com/in/your-handle}{linkedin.com/in/your-handle}
  }
\end{center}

\cvsection{教育背景}
\cvneedspace{3\baselineskip}
\textbf{学校名称} \hfill 2023 -- 2025 \\
硕士，专业名称

\cvneedspace{3\baselineskip}
\textbf{学校名称} \hfill 2019 -- 2023 \\
学士，专业名称

\cvsection{经历}
\cvneedspace{4\baselineskip}
\noindent\textbf{公司或实验室名称 \quad 岗位名称} \hfill 2024 -- 2025

{\small 用一段紧凑文字概括：你解决了什么问题，用了什么技术方案，交付了什么结果；仅在来源支持时使用量化结果。}

\vspace{3pt}
\cvneedspace{4\baselineskip}
\noindent\textbf{团队或机构名称 \quad 科研实习生 / 工程师 / 分析师} \hfill 2023 -- 2024

{\small 用一段紧凑文字概括第二段代表经历，强调职责边界、方法和结果。}

\cvsection{项目经历}
\cvneedspace{4\baselineskip}
\noindent\textbf{项目名称} \hfill 2025

{\small 用一段紧凑文字说明项目目标、系统结构、个人贡献和结果。}

\vspace{3pt}
\cvneedspace{4\baselineskip}
\noindent\textbf{项目名称} \hfill 2024

{\small 用一段更短的文字补充第二个代表项目。}

\cvsection{技能}
\textbf{机器学习 / AI：} 按目标岗位填写可自证技能。\\
\textbf{工程能力：} 按目标岗位填写语言、系统与工具。\\
\textbf{语言 / 其他：} 仅填写真实且相关的能力。

\end{document}
